\fancyhead[L]{\small MAGNETISCHE BEOBACHTUNGEN.}
\SourceLeaf{523}\pdfbookmark[0]{Nordlicht vom 7. Januar 1831}{nordlicht_1831_januar}
\begin{center}
\hrule\vspace{2mm}
Göttingische gelehrte Anzeigen. 1831. Februar 28.
\vspace{2mm}\hrule\vspace{5mm}
\end{center}
Herr Prof.\ \textsc{Gerling} in Marburg hat der Königl.\ Societät eine Notiz über seine Wahrnehmung
\begin{center}\itshape des am 7. Januar d.\ J. gesehenen Nordlichts\end{center}
\noindent vorgelegt, welche zwar im Allgemeinen mit dem, was von andern Orten her bereits bekannt geworden ist, übereinstimmt, aber daneben noch einen, besonderer Aufmerksamkeit werthen, und wie es scheint bisher noch nicht hinlänglich gewürdigten Umstand berührt, daher wir hier einen Auszug aus derselben mittheilen.

Das Phänomen war in Marburg schon von 6 Uhr an gesehen. Herr \textsc{Gerling} erhielt aber erst um 8 Uhr eine Benachrichtigung davon, und damals war am ganzen nordlichen Himmel, so tief herab wie die Aussicht aus den Fenstern seiner Wohnung reichte, gar nichts Ungewöhnliches zu erkennen. Allein gegen 9 Uhr zeigten sich wieder auffallende rothe Streifen am nordlichen Himmel, und Herr \textsc{Gerling} begab sich sogleich auf den eine freie Aussicht beherrschenden Schlossberg, um noch so viel thunlich von der Erscheinung wahrzunehmen.

Zuerst wurden in einer Ausdehnung von etwa 50—60 Grad zwischen N.O. und N.W. blos rothe Streifen und Flecken am Himmel bemerklich, welche sich ohne vollständige Continuität in dem angegebenen Bogen im Azimuth und im Mittel etwa bis zu 45 Grad Höhe erstreckten. In der Mitte jenes Azimuthalbo\SourceLeaf{524}gens um den Meridian herum und nach einer Schätzung etwa in 30—40 Grad Azimuthalausdehnung zeigten sich schwarze Flecke am sonst heitern Himmel, dem Ansehn nach mit nichts anderm als schwarzen Wölkchen zu vergleichen. Diese Flecke vermehrten sich allmählich, und bildeten endlich zusammenlaufend das dunkle Segment, welches nach allen Beschreibungen bei dem Nordlicht characteristisch zu sein scheint, indem zu gleicher Zeit die ersterwähnten rothen Flecken an Intensität zunahmen, und sich strahlenförmig gegen das schwarze Segment gruppirten, von welchem aus zwischen den rothen Strahlen dann auch weisse und gelbliche erschienen, die ohne auffallend plötzliches Fortschiessen sich auf etwa 50 Grad in der Höhe erstrecken mochten.

‘So weit, fährt Herr \textsc{Gerling} fort, scheint diese Beobachtung mit dem, was andere Beobachter zu gleicher Zeit und bei frühern Nordlichtern gesehen haben, ganz übereinzustimmen, und würde also kaum eine Erwähnung verdienen, wenn nicht ein Umstand dabei mir aufgefallen wäre, welcher meines Wissens weder bei Gelegenheit dieses jetzigen Nordlichts, noch, so viel ich habe auffinden können, sonst zur Sprache gekommen ist. Nemlich, nicht bloss die Sterne des Schwans, über welchen die weissen und rothen Strahlen mit ihrer grossen Intensität hinweggingen, sondern auch der Stern $\alpha$ in der Leyer, \textit{welcher tief im schwarzen Segment stand}, verloren an Sichtbarkeit und scheinbarer Helligkeit augenfällig gar nichts. Diese Thatsache scheint über die räthselhafte Frage, welche Bewandtniss es mit dem dunkeln Segment eigentlich habe, wenigstens das negative Resultat zu geben, dass es \textit{keine gewöhnliche Wolke} ist, weil solche für das Sternlicht nicht permeabel sein könnte.’

Schon bei dem Nordlicht vom 22. October 1804 bemerkte \textsc{Wrede}, allein ohne diesen Grund beizufügen, dass man das dunkle Segment \textit{unrichtig} eine Wolke nenne, während \textsc{Gilbert} den Ausdruck in Schutz nimmt, und hinzusetzt, er habe im dunkeln Segment nichts bemerkt, was ihn hätte auf den Gedanken bringen können, dass er dort etwas anderes als eine dunkle Wolke sähe. Auch die Meinung \textsc{Mayers} im Handbuch der physischen Astronomie, dass die dichtere mit Dünsten erfüllte Luft des Horizonts hinlänglich sei, das dunkle Segment zu erklären, scheint sich mit der von Herrn \textsc{Gerling} bemerkten Thatsache nicht vereinigen zu lassen.

Herr \textsc{Gerling} fügt noch bei, dass in den frühern Stunden, wo das in seiner Ausdehnung veränderliche Segment sich sehr hoch erstreckte, ein glaubwür\SourceLeaf{525}diger Zeuge den Stern $\alpha$ Leyer \textit{in dem Segmente} so hell wie zu irgend einer andern Zeit glänzen gesehen, und ein anderer, zu einer Zeit, wo das dunkle Segment sich noch nicht bis zu jenem Sterne erstreckte, \textit{andere Sterne in dem Segment} erblickt habe.

Herr \textsc{Gerling} hat noch einen Auszug aus seinem meteorologischen Journal vom 5.—9. Januar beigefügt, welcher jedoch ausser einem dreiviertel Zoll betragenden Steigen des Barometers vom 6. Januar Nachmittags bis 7. Januar Abends nichts auffallendes darbietet. Der Wind ging am 7. Januar aus Norden.

Die hier in Göttingen von Herrn Prof.\ \textsc{Harding} an diesem Nordlichte gemachten Wahrnehmungen stimmen im Wesentlichen mit den von andern Orten bekannt gewordenen überein, doch verdient der Umstand erwähnt zu werden, dass während der Dauer des Phänomens die Magnetnadel um etwa dreiviertel Grad von ihrer gewöhnlichen Stellung nach Norden ging, und am andern Morgen wieder auf dieselbe zurückgekommen war.

\bigskip
\clearpage
\pdfbookmark[0]{Magnetisches Observatorium. 1834}{magnetisches_observatorium_1834}
\begin{center}
\hrule\vspace{2mm}
Göttingische gelehrte Anzeigen. 1834. August 9.
\vspace{2mm}\hrule\vspace{5mm}
\end{center}
Wir verdanken der Huld unserer Regierung ein neues, einem wichtigen Theile der Naturwissenschaften gewidmetes Institut, ein eignes
\begin{center}\itshape für die magnetischen Beobachtungen und Messungen\end{center}
\noindent errichtetes Observatorium.

Obgleich der Bau desselben bereits im vorigen Herbst, und die innere Einrichtung seit Anfang dieses Jahrs so weit vollendet ist, dass seit den ersten Monaten tägliche Beobachtungen angestellt werden konnten, so haben wir doch bisher Anstand genommen, in diesen Blättern einen Bericht davon zu geben, weil wir erst einige Resultate der Beobachtungen damit verbinden zu können gewünscht haben. Die nach neuen Principien construirten magnetischen Apparate, welche im Jahre 1832 in der hiesigen Sternwarte aufgestellt sind, haben wir bereits früher in diesen Blättern [Anzeige d.\ \textit{Intensitas v.\,m.}] ausführlich beschrieben, und die damit
\SourceLeaf{526}erreichbare Schärfe ist aus dem dort Angeführten hinreichend ersichtlich: allein um diese Schärfe \textit{ganz} zu erreichen, war eine Ausführung in grösserm Maassstabe, und um den Resultaten eine vollkommene Reinheit von fremden Einflüssen zu verschaffen, war ein besonderes eisenfreies Gebäude unumgänglich nöthig.

Das magnetische Observatorium, auf einem freien Platze, etwa hundert Schritt westlich von der Sternwarte errichtet, ist ein genau orientirtes längliches Viereck von 32 Par.\ Fuss Länge und 15 Fuss Breite, mit zwei Vorsprüngen an den längeren Seiten; der westliche Vorsprung bildet den Eingang, und dient zugleich bei gewissen Beobachtungen als Erweiterung des Hauptsaals; der östliche Vorsprung, vom Hauptsaal ganz geschieden, dient zum Aufenthalt des Nachtwächters der Sternwarte. Im ganzen Gebäude ist ohne Ausnahme alles, wozu sonst Eisen verwandt wird, Schlösser, Thürangeln, Fensterbeschläge, Nägel u.s.w.\ von Kupfer. Für Abhaltung alles Luftzuges ist nach Möglichkeit gesorgt. Die Höhe des Saals ist etwas über 10 Fuss.

Der magnetische Apparat stimmt im Wesentlichen mit den oben erwähnten überein, daher wir uns darauf einschränken, nur die Verschiedenheiten anzugeben. Der Magnetstab ist aus Uslarschem Gussstahl, welcher sich zu magnetischen Versuchen vortrefflich qualificirt; es wird von Zeit zu Zeit mit verschiedenen Stäben gewechselt, die alle nahe gleiche Grösse haben, nemlich eine Länge von 610, Breite von 37, Dicke von 10 Millimetern; das Gewicht gegen vier Pfund. Der Spiegel ist 75 Millimeter breit und 50 hoch. Aufgehängt ist der Stab von der Mitte der Decke des Saals an einem 200fachen 7 Fuss langen ungedrehten Seidenfaden; der Torsionskreis ist aber nicht wie früher am obern Ende des Fadens, sondern am untern, und mit dem Schiffchen, welches den Stab trägt, drehbar verbunden. Seidene Aufhängungsfäden haben vor metallenen, wie bereits in der Abhandlung des Hofr.\ \textsc{Gauss} (\textit{Intensitas vis magneticae terrestris} Art.\ 9.) bemerkt ist, den grossen Vorzug, dass ihre Torsionskraft sehr klein ist; bei dem gegenwärtigen Tragfaden ist diese nur der Neunhundertste Theil der horizontalen Directionskraft des Magnetstabes, während die Torsionskraft eines Metallfadens von gleichem Tragvermögen etwa zehnmal stärker sein würde. Dagegen haben Seidenfäden, besonders wenn ihr Tragvermögen das an ihnen hängende Gewicht nicht weit übersteigt, die Inconvenienz, sich in den ersten Wochen, oder bei bedeutend verstärkter Belastung, beträchtlich zu verlängern; inzwischen wird dieser Inconvenienz hier durch den sinnreichen von Herrn Prof.\ \textsc{Weber} angege\SourceLeaf{527}benen an der Decke befindlichen Aufhängungsapparat abgeholfen, womit der Faden leicht, so viel nöthig, wieder aufgewunden werden kann, ohne seinen Platz zu verändern; zugleich aber kann dieser Apparat eben so leicht an der Decke verschoben werden, wenn im Lauf der Zeit die Veränderung der magnetischen Declination dies nöthig machen wird. Der Theodolith steht bisher auf einem sehr solide gearbeiteten hölzernen Stativ über einem besondern steinernen Fundament, und von dem Platze desselben ist durch das nördliche Fenster einer der Stadtthürme sichtbar, dessen Azimuth auf das genaueste bestimmt ist. Als Berichtigungsmarke für die unverrückte Stellung des Theodolithen dient blos ein zarter verticaler Strich an der gegenüberstehenden nördlichen Wand. Zum gewöhnlichen Gebrauch dient eine in Millimeter getheilte Scale von 4 Fuss Länge; für einige Beobachtungen wird dieselbe mit einer zwei Meter langen vertauscht. Der Werth eines Scalentheils ist $21''3$. Für nächtliche Beobachtungen wurde bisher die Scale mit starken Wachskerzen beleuchtet; in Zukunft werden dazu \textsc{Argand}sche Lampen gebraucht werden.

Eine der Hauptanwendungen des Apparats besteht nun in der scharfen Bestimmung der magnetischen Declination und ihrer Veränderung in verschiedenen Tagesstunden, Monaten und Jahren. Alle Tage wird die Aufzeichnung zweimal zu bestimmten Stunden gemacht: man hat dazu die Vormittagsstunde 8 Uhr, und die Nachmittagsstunde 1 Uhr gewählt, mit welchen Zeiten bei regelmässigem Verlauf der täglichen Variationen die kleinste und die grösste Declination, wenigstens in den ersten Monaten des Jahrs, ungefähr zusammenfallen: dieser Aufzeichnung allemal genau bei derselben Uhrzeit hat man, aus wichtigen hier nicht weiter auszuführenden Gründen, vor dem jedesmaligen Abwarten des Minimum und Maximum unbedingt den Vorzug geben müssen. Diese Aufzeichnungen haben zwar schon seit dem 1. Januar den Anfang genommen: allein da zuerst ein schwächerer Aufhängungsfaden angewandt war, dessen allmähliche Verlängerungen eine öftere Aufwindung nöthig machten, wobei nicht unbeträchtliche, Anfangs nicht genug beachtete Veränderungen des Nullpunkts der Torsion eingetreten sind, so hat man die ersten drittehalb Monate lieber ausgeschlossen. Die seitdem erhaltenen Mittelwerthe für die westliche Declination der Magnetnadel sind folgende gewesen:
\SourceLeaf{528}
\begin{center}
\begin{tabular}{l|rrr|rrr}
 & \multicolumn{3}{c|}{8 Uhr. Vorm.} & \multicolumn{3}{c}{1 Uhr Nachm.}\\\hline
März, zweite Hälfte & $18^0$ & $38'$ & $16''0$ & $18^0$ & $46'$ & $40''4$\\
April & & $36$ & $6.9$ & & $47$ & $3.8$\\
Mai & & $36$ & $28.2$ & & $47$ & $15.4$\\
Junius & & $37$ & $40.7$ & & $47$ & $59.5$\\
Julius & & $37$ & $57.5$ & & $48$ & $19.0$
\end{tabular}
\end{center}

Ferner werden an gewissen bestimmten Tagen im Jahre 44 Stunden hindurch ununterbrochen in kurzen Zeitfristen die Veränderungen der Declination beobachtet. Man hat dazu dieselben bereits vor mehreren Jahren durch Herrn von \textsc{Humboldt} festgesetzten Tage gewählt, an welchen nach Verabredung schon an vielen zum Theil sehr entlegenen Plätzen ähnliche Aufzeichnungen mit \textsc{Gambey}schen Apparaten gemacht werden. Bis jetzt sind hier diese Beobachtungen dreimal angestellt, nemlich den 20. 21. März; 4. 5. Mai; 21. 22. Junius, und es haben daran Theil genommen ausser dem Hofr.\ \textsc{Gauss} die Herren Prof.\ \textsc{Weber}, Prof.\ \textsc{Ulrich}, Dr.\ \textsc{Weber}, Dr.\ \textsc{Goldschmidt}, Dr.\ \textsc{Listing}, \textsc{Sartorius}, \textsc{Deahna} und \textsc{Wilh.\ Gauss}. Der Zweck dieser Beobachtungen ist, theils den regelmässigen Verlauf nach und nach immer vollständiger kennen zu lernen, theils die Bewandtniss, welche es mit den so häufig dazwischen kommenden, zuweilen, besonders bei Nordlichtern, ungemein beträchtlichen ausserordentlichen Anomalien hat, durch Vergleichung der gleichzeitigen Beobachtungen an verschiedenen Orten zu erforschen. Die Aufzeichnungen geschahen hier, im März von 20 zu 20 Minuten, und zum Theil in halb so grossen Zwischenzeiten; im Mai von 10 zu 10 Minuten und zum Theil in doppelt engen Grenzen; im Junius durchgehends von 5 zu 5 Minuten. Anomalien wurden hier bemerkt, ein Paar auffallend grosse in der Nacht vom 20. zum 21. März; sehr bedeutende und zahlreiche in den Nächten vom 4. und 5. Mai; und einige zwar nicht grosse aber doch bestimmt hervortretende am 21. Junius, während den ganzen 22. Junius der Verlauf überaus regelmässig war. Von denjenigen correspondirenden Beobachtungen, welche, wie schon erwähnt, Herrn von \textsc{Humboldt} ihre Veranlassung verdanken, sind uns bisher keine bekannt geworden, als die Berliner vom 20. 21. März, welche jedoch nur von Stunde zu Stunde aufgezeichnet waren, und daher keine besondere Resultate geben konnten, obwohl sie doch eine Andeutung der in Göttingen bemerkten und verfolgten Anomalien enthielten. Dagegen wurden von Herrn \textsc{Sartorius} mit einem zwar kleinern aber nach denselben Principien wie
\SourceLeaf{529}der hiesige construirten Apparate die correspondirenden Beobachtungen vom 4. und 5. Mai auf einem Gute in Baiern, einige Meilen südlich von Meiningen sehr vollständig angestellt, woraus eine wahrhaft bewundernswürdige Übereinstimmung mit den hier beobachteten grossen Anomalien, nach Zeit, Grösse und Wechsel derselben hervorgeht, so dass man in den graphischen Darstellungen die eine beinahe als eine Copie der andern mit allen barocken durch jene Anomalien hervortretenden Figuren ansehen möchte. Ein eben so schöner Erfolg hat sich am 21. und 22. Junius gezeigt, wo correspondirende Beobachtungen in Berlin zum ersten Male mit einem dem hiesigen ähnlichen obwohl kleinern Apparate von Herrn Prof.\ \textsc{Encke} unter Beistand von Herrn \textsc{Poggendorff}, \textsc{Mädler} und \textsc{Wolfers} angestellt wurden. Auch dort waren keine andere Anomalien, als die hier beobachteten, aber diese fast treu copirt, und eben dasselbe zeigten die von Herrn \textsc{Sartorius} dasmal in Frankfurt am Main gemachten Beobachtungen. Diese Resultate können bereits als eine schöne Frucht der verabredeten Beobachtungen angesehen werden, da daraus auf das klarste hervorgeht, dass kleinere und grössere Anomalien der Magnetnadel, die zuweilen in ziemlich kurzen Fristen wechseln, nicht locale, sondern kräftige, weithin wirkende Ursachen haben müssen, was man in Beziehung auf sehr grosse mit Nordlichtern in Verbindung stehenden Unregelmässigkeiten auch schon früher bemerkt hatte. So wie in Zukunft die Theilnahme an diesen verabredeten Beobachtungen mit den eben so scharfen als bequemen Apparaten sich immer weiter ausbreiten wird, wozu schon die schönsten Aussichten vorhanden sind, wird es nicht fehlen, dass wir über diese höchst merkwürdigen und räthselhaften Erscheinungen umfassende Aufklärungen erhalten.

Übrigens werden hier solche Beobachtungen auch ausser den bestimmten Zeiten häufig gemacht, wobei zuweilen ganz auffallende Anomalien vorgekommen sind. So nahm z.\ B. am 14. Januar Abends zwischen 8 und 9 Uhr die Declination innerhalb Einer Viertelstunde um 13 Minuten mit grösster Regelmässigkeit ab, und kehrte dann allmählich auf ihren vorigen Stand zurück. Dergleichen Wahrnehmungen können indess keine weitere Resultate geben, da ohne Verabredung correspondirende Beobachtungen höchst selten zu erwarten sind.

Von Zeit zu Zeit wird in dem hiesigen magnetischen Observatorium auch die Bestimmung der absoluten Intensität des Erdmagnetismus wiederholt werden. Da, um diese Operation mit grösster Schärfe auszuführen, erst verschie\SourceLeaf{530}dene Vorkehrungen getroffen werden mussten, so hat sie das erste Mal erst im Julius gemacht werden können. Drei Bestimmungen mit verschiedenen Stäben gaben
\begin{center}
\begin{tabular}{rllllllr}
17. Julius & . & . & . & . & . & & $1.7743$\\
20. \enspace „ & . & . & . & . & . & & $1.7740$\\
21. \enspace „ & . & . & . & . & . & & $1.7761$
\end{tabular}
\end{center}
als Werth der horizontalen Kraft, wobei, wie bei den frühern Bestimmungen mit kleinern Stäben, deren geringe Verschiedenheit von den gegenwärtigen man mit Vergnügen bemerken wird, die Zeitsecunde, das Millimeter und das Milligramm als Einheiten zum Grunde liegen.

Eben so, wie mit dem frühern in der Sternwarte aufgestellten Apparate, hat man nun auch mit dem gegenwärtigen im magnetischen Observatorium Vorrichtungen zu electro magnetischen Versuchen und Messungen verbunden. Der aufgehängte Magnetstab ist von einem aus 200 Umwindungen bestehenden Multiplicator umgeben, dessen Construction die Anwendung von nichtbesponnenem Draht erlaubte: die Drahtlänge beträgt 1100 Fuss. Mit Hülfe eines sehr einfach construirten Commutators kann der Beobachter, ohne sein Auge vom Fernrohr zu entfernen, jeden Augenblick die Richtung des galvanischen Stroms umkehren, oder den Strom ganz unterbrechen.

Wir können hiebei eine mit den beschriebenen Einrichtungen in genauer Verbindung stehende grossartige und bisher in ihrer Art einzige Anlage nicht unerwähnt lassen, die wir unserm Herrn Prof.\ \textsc{Weber} verdanken. Dieser hatte bereits im vorigen Jahre von dem physikalischen Cabinet aus über die Häuser der Stadt hin bis zur Sternwarte eine doppelte Drahtverbindung geführt, welche gegenwärtig von der Sternwarte bis zum magnetischen Observatorium fortgesetzt ist. Dadurch bildet sich eine grosse galvanische Kette, worin der galvanische Strom, die an beiden Endpunkten befindlichen Multiplicatoren mitgerechnet, eine Drahtlänge von fast neuntausend Fuss zu durchlaufen hat. Der Draht der Kette ist grösstentheils Kupferdraht von der im Handel mit 3 bezeichneten Nummer, wovon eine Länge von einem Meter acht Gramm wiegt; der Draht des Multiplicators im magnetischen Observatorium ist übersilberter Kupferdraht Nr.\ 14, wovon auf ein Gramm 2,6 Meter kommen. Diese Anlage ist ganz dazu geeignet, zu einer Menge der interessantesten Versuche Gelegenheit zu geben. Man be\SourceLeaf{531}merkt nicht ohne Bewunderung, wie ein einziges Plattenpaar am andern Ende hineingebracht, augenblicklich dem Magnetstabe eine Bewegung ertheilt, die zu einem Ausschlage von weit über tausend Scalentheilen ansteigt; noch auffallender aber findet man wenigstens anfangs, dass ein Plattenpaar von sehr geringer Grösse, z.\ B. Einen Zoll im Durchmesser, und unter Anwendung von blossem Brunnen- oder selbst destillirten Wasser eine nicht viel kleinere Wirkung hervorbringt, als ein sehr grosses Plattenpaar mit starker Säure. Und doch ist dieser Umstand bei näherer Überlegung ganz in der Ordnung, und dient nur zu neuer Bestätigung der schönen zuerst von \textsc{Ohm} aufgestellten Theorie. Bei Vermehrung der Anzahl der Plattenpaare wächst hingegen die Wirkung, und zwar dieser beinahe proportional. Die Leichtigkeit und Sicherheit, womit man durch den Commutator die Richtung des Stroms und die davon abhängige Bewegung der Nadel beherrscht, hatte schon im vorigen Jahre Versuche einer Anwendung zu telegraphischen Signalisirungen veranlasst, die auch mit ganzen Wörtern und kleinen Phrasen auf das vollkommenste gelangen. Es leidet keinen Zweifel, dass es möglich sein würde, auf ähnliche Weise eine unmittelbare telegraphische Verbindung zwischen zweien eine beträchtliche Anzahl von Meilen von einander entfernten Örtern einzurichten: allein es kann natürlich hier nicht der Ort sein, Ideen über diesen Gegenstand weiter zu entwickeln.

\bigskip
\pdfbookmark[0]{Beobachtungen der magnetischen Variation. 1834}{variation_1834_october}
\begin{center}\itshape
Beobachtungen der magnetischen Variation in Göttingen und Leipzig\\
am 1. und 2. October 1834.
\end{center}
\begin{center}
\hrule\vspace{2mm}
\textsc{Poggendorff}. Annalen der Physik und Chemie. 1834. Bd.\ 33.
\vspace{2mm}\hrule\vspace{5mm}
\end{center}
Die in meinem Aufsatze über das hiesige magnetische Observatorium erwähnten Beobachtungen der magnetischen Variation an den verabredeten Tagen sind seitdem hier noch zwei Mal angestellt, am 6. und 7. August, und am 23. und 24. September. Im ersten Termin kamen recht starke und merkwürdige Anomalien vor, und es ist daher um so mehr zu bedauern, dass zufällige Ursachen die Anstellung correspondirender Beobachtungen an andern Orten gehindert haben. Die September-Beobachtungen sind hingegen ganz vollständig auch
